\chapter{Introduction}

Reliable radio coverage prediction is a central requirement for future 6G
systems in which UAVs may act as aerial base stations, relays, or temporary
coverage nodes. Ray tracing can produce detailed channel knowledge maps, but it
is too expensive to repeat for every transmitter height, city morphology or
deployment hypothesis. This thesis studies whether a learned surrogate can
predict dense channel maps directly from urban geometry while preserving enough
physical structure to generalize to cities not seen during training.

The project is formulated as dense prediction on \(513\times513\) city maps.
Inputs include building-height topology, geometric LoS/NLoS masks and the UAV
antenna height. The outputs are path loss, delay spread and angular spread. The
final line of work combines train-only calibrated priors with a residual
deep-learning model, so the neural network corrects a strong physical/statistical
anchor instead of learning the entire propagation field from pixels alone.

\section{Goals and research questions}

The original success targets were kept throughout the project:
path-loss RMSE below \SI{5}{dB}, delay-spread RMSE below \SI{50}{ns},
angular-spread RMSE below \(20^\circ\), separate LoS/NLoS reporting, and a clear
runtime gain over ray tracing. The final model meets those targets on unseen
test cities, reaching \SI{1.6519}{dB}, \SI{26.5570}{ns} and \(11.3854^\circ\)
RMSE for PL, DS and AS, respectively.

The thesis answers three questions:
\begin{itemize}
  \item \textbf{RQ1.} Can one dense surrogate predict PL, DS and AS under strict
        city holdout? Yes, within the fixed CKM simulation contract, using a
        shared prior-anchored residual GMM-head model.
  \item \textbf{RQ2.} What should be handled by priors and what by deep
        learning? Stable radial, visibility and height structure is best handled
        by frozen calibrated priors; the network is most useful as a bounded
        residual corrector.
  \item \textbf{RQ3.} Can continuous UAV height be handled by one model? Yes,
        provided height is explicit in the priors, FiLM conditioning and
        evaluation; low-altitude and NLoS-heavy cases remain the hard regimes.
\end{itemize}

The main contributions are a reproducible city-holdout CKM evaluation contract,
a height-aware calibrated path-loss prior, log-domain calibrated DS/AS priors,
and a joint residual GMM-head model that improves those frozen anchors.

\section{Dataset and evaluation contract}

The project uses the CKM HDF5 dataset \texttt{CKM\_Dataset\_270326.h5}, with 66
real-city scenarios and 16180 samples. Each sample is a \(513\times513\) raster
at \SI{1}{m}/pixel. The UAV transmitter is horizontally centred, but its
antenna height varies continuously from about \SI{12}{m} to \SI{478}{m}. Every
non-building pixel is treated as a receiver at \(h_{\mathrm{rx}}=\SI{1.5}{m}\);
building pixels remain as propagation context but are excluded from losses and
metrics.

\begin{table}[htbp]
\centering
\caption{Main CKM fields used by the thesis.}
\label{tab:dataset_fields_intro}
\scriptsize
\begin{tabularx}{\textwidth}{lllX}
\toprule
Field & Shape & Type & Role \\
\midrule
\texttt{topology\_map} & \(513^2\) & float32 & Building height and ground mask source. \\
\texttt{los\_mask} & \(513^2\) & uint8 & Geometric LoS/NLoS visibility. \\
\texttt{path\_loss} & \(513^2\) & uint8 & PL in dB, approximately \SI{1}{dB}-quantized. \\
\texttt{delay\_spread} & \(513^2\) & uint16 & DS target, decoded/evaluated in ns. \\
\texttt{angular\_spread} & \(513^2\) & uint8 & AS target, decoded/evaluated in degrees. \\
\texttt{uav\_height} & scalar & float32 & UAV height for priors and FiLM conditioning. \\
\bottomrule
\end{tabularx}
\end{table}

The key methodological choice is the city split. Complete cities are reserved
for validation and test, so the model must transfer to unseen layouts instead of
memorizing similar tiles from the same city. This contract is used for all final
model-versus-prior comparisons.

\section{Development route and thesis structure}

The initial plan was a standard image-to-image pipeline: literature review,
dataset preparation, U-Net/cGAN baselines, model improvement and evaluation. The
actual route became more diagnostic. Early CNNs produced plausible maps but did
not control physical RMSE; the ground-mask reset changed the metric; PMHHNet and
topology experts improved LoS but not NLoS tails; distribution-first GMM heads
showed that the target distributions were multi-modal or spike-plus-tail; and
the final solution returned to strong frozen priors plus bounded neural
residuals.

Chapter~2 positions this work against dense radio-map prediction, hybrid
propagation models and distribution-aware heads. Chapter~3 gives the compact
methodology, including the detailed Try 78/79 priors and Try 80 residual GMM
head. Chapter~4 reports final unseen-city results, LoS/NLoS diagnostics,
height/topology audits, SOA context and runtime. The remaining chapters cover
sustainability, conclusions and compact appendices; longer attempt logs,
manual-style generator documentation and full qualitative galleries remain in
the source archive.
